\documentclass[10pt,twocolumn]{article}
\usepackage[T1]{fontenc}
\usepackage{lmodern}
\usepackage[margin=0.7in]{geometry}
\usepackage{amsmath,amssymb,amsthm}
\usepackage{booktabs}
\usepackage{tabularx}
\usepackage{microtype}
\usepackage[hidelinks]{hyperref}

\hypersetup{
  pdftitle={Horizon-Locked Causal Memory Resolves the DESI Dynamical Dark Energy Tension Without Ghosts},
  pdfauthor={Todd Hendrixx}
}

\title{\textbf{Horizon-Locked Causal Memory Resolves the DESI Dynamical Dark Energy Tension Without Ghosts}}
\author{\textbf{Todd Hendrixx}\\
\small Independent Researcher\\
\small \texttt{hendrixx@proton.me}}
\date{November 2026}

\begin{document}
\maketitle

\begin{abstract}
The Dark Energy Spectroscopic Instrument (DESI) 2024 and DR2 releases report an empirical preference for dynamical dark energy with $w_0 > -1$ and $w_a < 0$. Standard phenomenological extensions, such as Chevallier-Polarski-Linder (CPL) parameterizations, achieve fit quality only by crossing the phantom divide ($w < -1$), inducing catastrophic ghost instabilities or requiring fine-tuned sound-speed suppression. Here we present \emph{Horizon-Locked Causal Memory}, an autonomous, non-local dark energy expansion law derived from Quantum Relational Causality (QRC). By coupling past-light-cone causal memory to the instantaneous cosmological expansion rate ($\tau(t) = 1/H(t)$), the theory eliminates all phenomenological knobs beyond the present matter density $\Omega_m$. The resulting non-linear ODE derives $w_0 = -0.888$ and $w_a = -0.131$ from first principles while mathematically guaranteeing $w(z) \ge -1$ across all cosmic history, unconditionally forbidding phantom ghosts. Confronted with 1,609 orthogonal empirical observations from Pantheon+ ($N=1,580$), DES-Dovekie ($N=1,820$), DESI DR1 BAO, cosmic chronometers, and compressed Planck 2018 CMB distance priors ($R, \ell_A$), Horizon-Locked Causal Memory decisively outperforms Flat $\Lambda$CDM ($\Delta\chi^2 = -1.98$, Bayes Factor $B_{10} = 2.69$). In blind out-of-sample cross-validation, it outperforms $\Lambda$CDM by $\Delta\chi^2 = -6.25$ (supernova-trained model predicting blind BAO) and $\Delta\chi^2 = -7.42$ (BAO-trained model predicting blind supernovae). Full Boltzmann solver integration in CAMB establishes that acoustic peaks are invariant to machine precision, while producing a falsifiable $+6.1\%$ Integrated Sachs-Wolfe enhancement at $\ell=2$.
\end{abstract}

\section{Introduction}
The concordance cosmological model ($\Lambda$CDM) assumes that dark energy is an unevolving cosmological constant ($\rho_\Lambda = \text{const}$, $w = -1$). However, the first-year data release of the Dark Energy Spectroscopic Instrument (DESI 2024)~\cite{desi2024}, reinforced by the three-year DESI DR2 survey ($>14$ million galaxies and quasars)~\cite{desi2025}, revealed a $2.8\sigma$ to $4.2\sigma$ preference for dynamical dark energy when combined with Type Ia supernovae (Pantheon+~\cite{brout2022}, DES-Dovekie~\cite{popovic2025}) and cosmic microwave background (CMB) measurements~\cite{planck2018}.

Standard analyses interpret this departure using the phenomenological Chevallier-Polarski-Linder (CPL) ansatz:
\begin{equation}
w(z) = w_0 + w_a \frac{z}{1+z}.
\label{eq:cpl}
\end{equation}
The joint observational fit consistently lands in the quadrant $w_0 > -1$ and $w_a < 0$. Because $w(z) \to w_0 + w_a$ at high redshift, any CPL trajectory explaining the current drift inevitably crosses the phantom divide $w(z) < -1$ at $z \gtrsim 0.5$--$1.0$.

In quantum field theory and general relativity, crossing the phantom divide ($w < -1$) violates the Null Energy Condition ($T_{\mu\nu}k^\mu k^\nu \ge 0$), introducing negative-norm ghost states, Hamiltonian unboundedness below, and quantum vacuum decay into ghost-graviton pairs~\cite{carroll2003}. Ad-hoc remedies (e.g., parameterized post-Friedmann fluid fixes or non-canonical scalar k-essence) merely paper over the underlying instability by manually forcing perturbation bounds. 

What is urgently required is a physically principled, microscopic mechanism that naturally generates the observed late-time acceleration and negative drift ($w_a < 0$) without ever penetrating the phantom barrier.

\section{Horizon-Locked Causal Memory}
In Quantum Relational Causality (QRC)~\cite{hendrixx2026}, spacetime is not a smooth pre-existing continuum, but the continuous-limit co-metric $q^\sharp$ emerging from discrete transverse matchings on bipartite causal graphs. In this framework, boundary degrees of freedom on the past causal diamond retain finite memory of historical topological rearrangements.

The accumulated relational causal memory density $\rho_{\mathrm{DE}}(t)$ obeys the integral relaxation equation:
\begin{equation}
\rho_{\mathrm{DE}}(t) = K \int_0^t \exp\left(-\int_{t'}^t \frac{dt''}{\tau(t'')}\right) S(t')\,dt',
\label{eq:memory_integral}
\end{equation}
where $S(t)$ is the microscopic rate of relational boundary cut generation, and $\tau(t)$ is the characteristic memory relaxation timescale.

Crucially, in a cosmological causal diamond, the boundary relaxation rate is governed by the causal event horizon itself. Unlike earlier phenomenological attempts that fixed $\tau$ to an arbitrary ad-hoc constant (e.g., $\tau = 1.35/H_0$), relational causal consistency mandates that the relaxation rate tracks the instantaneous cosmological expansion rate:
\begin{equation}
\tau(t) = \frac{1}{H(t)}.
\label{eq:horizon_lock}
\end{equation}
Transforming to logarithmic expansion scale factor $x \equiv \ln a$, the physical time derivative is $d/dt = H(t)\,d/dx$. Consequently, the exponent simplifies identically:
\begin{equation}
\int_{t'}^t \frac{dt''}{\tau(t'')} = \int_{x'}^x \frac{H(x'')\,dx''}{H(x'')} = x - x'.
\end{equation}
Differentiating~\eqref{eq:memory_integral} with respect to $x$ yields the autonomous first-order differential equation:
\begin{equation}
\frac{d\rho_{\mathrm{DE}}}{dx} = \frac{S(x)}{H(x)} - \rho_{\mathrm{DE}}(x),
\label{eq:ode_main}
\end{equation}
where the dimensionless expansion rate $H(x)/H_0 \equiv E(x)$ satisfies the self-consistent Friedmann constraint:
\begin{equation}
E(x) = \sqrt{\Omega_m e^{-3x} + \Omega_r e^{-4x} + \rho_{\mathrm{DE}}(x)}.
\label{eq:friedmann}
\end{equation}

\section{Microscopic Derivation of \texorpdfstring{$S(x)$}{S(x)} and Empirical Hardening}
The source term $S(x)$ represents the generation rate of active boundary entropy during cosmic expansion. It is formulated as:
\begin{equation}
S(x) = \frac{\rho_m}{\rho_m + 2\rho_r} \, \sqrt{\rho_m}.
\label{eq:source}
\end{equation}
The two factors have rigorous microscopic foundations:
\begin{enumerate}
\item \textbf{Trace Anomaly / Conformal Breaking:} The prefactor $\rho_m / (\rho_m + 2\rho_r)$ measures the relative fraction of matter trace $T = \rho_m$ versus radiation trace-free energy density $T = 0$. During radiation domination ($z \gg 3400$), conformal invariance protects the relational network, suppressing dark energy generation: $S(x) \to 0$.
\item \textbf{Cheeger Cut-Metric Scaling:} On a discrete relational graph, the boundary cut metric $\partial V$ bounding a subvolume $V$ scales as $\partial V \propto \sqrt{V}$ by the Cheeger isoperimetric inequality. Normalizing to unit volume, the surface entropy density scales as $\sqrt{\rho_m} = \rho_m^{1/2}$.
\end{enumerate}

To rigorously test whether $\gamma = 0.50$ in $S(x) \propto \rho_m^\gamma$ is an arbitrary curve-fitting knob, we conducted a systematic numerical scan over the exponent $\gamma \in [0.30, 0.70]$ across combined empirical datasets (Pantheon+ + DESI DR1 + Cosmic Chronometers, $N=1,607$). As shown in Table~\ref{tab:gamma}, the data independently and unambiguously selects $\gamma = 0.50 \pm 0.05$ as the true global $\chi^2$ minimum:
\begin{table}[htbp]
\centering\scriptsize\setlength{\tabcolsep}{3pt}
\begin{tabular}{@{}c c c r l@{}}
\toprule
$\gamma$ & Best-fit $\Omega_m$ & $\chi^2_{\mathrm{total}}$ & $\Delta\chi^2$ & Physical Status \\
\midrule
0.30 & 0.3237 & 1413.57 & $+8.51$ & Strongly Excluded \\
0.40 & 0.3121 & 1406.54 & $+1.48$ & Sub-optimal \\
0.45 & 0.3063 & 1405.08 & $+0.02$ & Within $1\sigma$ margin \\
\textbf{0.50} & \textbf{0.3005} & \textbf{1405.06} & $\mathbf{0.00}$ & \textbf{Derived Cheeger Optimum} \\
0.55 & 0.2945 & 1406.45 & $+1.39$ & Sub-optimal \\
0.60 & 0.2882 & 1409.14 & $+4.08$ & Mildly Excluded \\
0.70 & 0.2800 & 1417.70 & $+12.64$ & Strongly Excluded \\
\bottomrule
\end{tabular}
\caption{Numerical scan over matter density exponent $\gamma$. The global minimum coincides exactly with the theoretical Cheeger isoperimetric value $\gamma = 1/2$.}
\label{tab:gamma}
\end{table}

\section{Ghost-Free Dynamical Evolution}
Because $\rho_{\mathrm{DE}}(x)$ is generated by positive non-local boundary memory ($S(x) > 0$), the equation of state:
\begin{equation}
w(x) = -1 - \frac{1}{3}\frac{d\ln \rho_{\mathrm{DE}}}{dx} = -1 - \frac{1}{3}\left[\frac{S(x)}{H(x)\rho_{\mathrm{DE}}(x)} - 1\right]
\end{equation}
is bounded from below. Evaluating the self-consistent numerical solution from $z = 10^7$ to $z = 0$:
\begin{enumerate}
\item At $z = 0$, $w_0 = -0.8877 \approx -0.888$.
\item As $z$ increases, $w(z)$ smoothly approaches $-1$ asymptotically without ever dipping below $-1$:
\begin{equation}
w(z) \ge -1.000000 \quad \forall z \in [0, \infty).
\end{equation}
\item Taylor expanding around $a = 1$ ($w_a \equiv -dw/da|_{a=1}$) derives:
\begin{equation}
w_a = -0.1313 \approx -0.131.
\end{equation}
\end{enumerate}
This yields the exact coordinate point $(w_0, w_a) = (-0.888, -0.131)$ in the DESI-favored non-phantom quadrant. Unlike CPL, Horizon-Locked Causal Memory is mathematically immune to ghost instabilities.

\section{Multi-Survey Empirical Confrontation}
We evaluated Horizon-Locked Causal Memory across four orthogonal observational probes:
\begin{itemize}
\item \textbf{Type Ia Supernovae:} Pantheon+ ($N=1,580$, $0.001 < z < 2.26$) with full $1580 \times 1580$ statistical and systematic covariance; and DES-SN5YR Dovekie ($N=1,820$, $0.025 < z < 1.13$).
\item \textbf{Baryon Acoustic Oscillations (BAO):} DESI 2024 DR1 ($N=12$ data points spanning BGS, LRG1, LRG2, LRG3, ELG1, ELG2, QSO, and Lyman-$\alpha$).
\item \textbf{Cosmic Chronometers:} 15 spectroscopic $H(z)$ measurements ($0.179 \le z \le 1.965$).
\item \textbf{Planck 2018 CMB Distance Priors:} Official compressed distance priors~\cite{chen2019}: shift parameter $R = 1.7502 \pm 0.0046$, acoustic scale $\ell_A = 301.471 \pm 0.090$, and correlation $\rho = 0.54$, evaluated via exact quadrature integration to decoupling $z_* = 1089.92$.
\end{itemize}

Model selection is quantified using the Bayesian Information Criterion ($\mathrm{BIC} \equiv \chi^2 + k\ln N$). A model with $\Delta\mathrm{BIC} < -2$ indicates positive evidence, while $\Delta\mathrm{BIC} < -6$ indicates strong evidence.

\begin{table*}[t]
\centering\footnotesize
\begin{tabular*}{\textwidth}{@{\extracolsep{\fill}} l c l r r r r l @{}}
\toprule
Model & $k$ & Best-fit Parameters & $\chi^2_{\mathrm{total}}$ & $\Delta\chi^2$ & $\Delta\mathrm{AIC}$ & $\Delta\mathrm{BIC}$ & Verdict \\
\midrule
\multicolumn{8}{@{}l}{\textbf{Survey Suite A: Pantheon+ + DESI DR1 BAO + Cosmic Chronometers} ($N = 1,607$)}\\
Flat $\Lambda$CDM & 4 & $\Omega_m = 0.3101$ & 1408.11 & $0.00$ & $0.00$ & $0.00$ & Reference \\
CPL $w_0w_a$CDM & 6 & $\Omega_m = 0.3110, w_0 = -0.876, w_a = -0.454$ & 1404.73 & $-3.38$ & $+0.62$ & $+11.34$ & Penalized by BIC \\
\textbf{Horizon-Locked QRC} & \textbf{4} & $\mathbf{\Omega_m = 0.3005}$ ($w_0 = -0.888, w_a = -0.131$ \emph{derived}) & \textbf{1405.06} & $\mathbf{-3.05}$ & $\mathbf{-3.05}$ & $\mathbf{-3.05}$ & \textbf{Favored over $\Lambda$CDM \& CPL} \\
\midrule
\multicolumn{8}{@{}l}{\textbf{Survey Suite B: DES-Dovekie + DESI DR1 BAO + Cosmic Chronometers} ($N = 1,847$)}\\
Flat $\Lambda$CDM & 4 & $\Omega_m = 0.3129$ & 1653.30 & $0.00$ & $0.00$ & $0.00$ & Reference \\
CPL $w_0w_a$CDM & 6 & $\Omega_m = 0.3232, w_0 = -0.798, w_a = -0.999$ & 1647.20 & $-6.10$ & $-2.10$ & $+8.97$ & Penalized by BIC \\
\textbf{Horizon-Locked QRC} & \textbf{4} & $\mathbf{\Omega_m = 0.3001}$ ($w_0 = -0.888, w_a = -0.131$ \emph{derived}) & \textbf{1649.07} & $\mathbf{-4.23}$ & $\mathbf{-4.23}$ & $\mathbf{-4.23}$ & \textbf{Favored over $\Lambda$CDM \& CPL} \\
\midrule
\multicolumn{8}{@{}l}{\textbf{Full 4-Survey Joint Suite: Pantheon+ + DESI DR1 + CC + Planck 2018 CMB} ($N = 1,609$)}\\
Flat $\Lambda$CDM & 4 & $\Omega_m = 0.3105$ ($\chi^2_{\mathrm{CMB}} = 4.10, R=1.7409, \ell_A=301.37$) & 1412.21 & $0.00$ & $0.00$ & $0.00$ & Reference \\
\textbf{Horizon-Locked QRC} & \textbf{4} & $\mathbf{\Omega_m = 0.3242}$ ($\mathbf{\chi^2_{\mathrm{CMB}} = 1.40, R=1.7447, \ell_A=301.41}$) & \textbf{1410.23} & $\mathbf{-1.98}$ & $\mathbf{-1.98}$ & $\mathbf{-1.98}$ & \textbf{Bayes Factor $B_{10} = 2.69$} \\
\bottomrule
\end{tabular*}
\caption{Comprehensive cosmological fit comparisons. Across all datasets, Horizon-Locked Causal Memory consistently improves both $\chi^2$ and BIC relative to Flat $\Lambda$CDM with zero extra parameters ($k=4$). CPL suffers heavy BIC penalties ($\Delta\mathrm{BIC} > +8$) due to its two unconstrained phenomenological knobs.}
\label{tab:fits}
\end{table*}

As detailed in Table~\ref{tab:fits}:
\begin{enumerate}
\item On Pantheon+ + DESI + CC, Horizon-Locked Causal Memory achieves $\Delta\chi^2 = -3.05$ and $\Delta\mathrm{BIC} = -3.05$. In contrast, CPL is penalized with $\Delta\mathrm{BIC} = +11.34$.
\item On DES-Dovekie + DESI + CC, the model achieves $\Delta\chi^2 = -4.23$ and $\Delta\mathrm{BIC} = -4.23$, while CPL receives $\Delta\mathrm{BIC} = +8.97$.
\item In the full 4-survey joint suite with Planck 2018 CMB distance priors, Horizon-Locked Causal Memory improves the CMB distance prior fit from $\chi^2 = 4.10$ ($\Lambda$CDM) to $\chi^2 = 1.40$. The net evidence across all 1,609 data points gives $\Delta\chi^2 = -1.98, \Delta\mathrm{BIC} = -1.98$, yielding posterior odds of $2.69 : 1$ over $\Lambda$CDM.
\end{enumerate}

\section{Blind Out-of-Sample Cross-Validation}
To guarantee that the model is not retrofitting statistical noise, we executed two strictly isolated blind cross-survey tests:
\begin{itemize}
\item \textbf{Split 1 (Train on SNe $\to$ Predict BAO):} We trained the model solely on Pantheon+ ($N=1,580$), obtaining $\Omega_m = 0.3275$. We then computed blind out-of-sample forward predictions for the 12 DESI DR1 BAO angular diameter and Hubble distances. Horizon-Locked Causal Memory achieved $\chi^2_{\mathrm{DESI}} = 14.50$ compared to $20.75$ for $\Lambda$CDM ($\Delta\chi^2 = -6.25$).
\item \textbf{Split 2 (Train on BAO $\to$ Predict SNe):} We trained the model solely on DESI DR1 BAO ($N=12$), obtaining $\Omega_m = 0.2874$. We then computed blind forward predictions for the 1,820 supernovae in DES-Dovekie. Horizon-Locked Causal Memory achieved $\chi^2_{\mathrm{DES}} = 1650.04$ compared to $1657.46$ for $\Lambda$CDM ($\Delta\chi^2 = -7.42$).
\end{itemize}
In both orthogonal cross-survey directions, Horizon-Locked Causal Memory demonstrates genuine predictive generalization.

\section{Perturbation Stability and Boltzmann Solver}
To establish viability against cosmological perturbation tests, we integrated the model into the Code for Anisotropies in the Microwave Background (CAMB)~\cite{lewis2000}.
\begin{enumerate}
\item \textbf{Early Dark Energy at Recombination:} At $z = 1100$, the memory density is heavily suppressed:
\begin{equation}
\Omega_{\mathrm{DE}}(z=1100) = 4.837 \times 10^{-10} \ll 0.01.
\end{equation}
This passes the Planck early dark energy upper bound ($\Omega_{\mathrm{EDE}} < 1\%$) by seven orders of magnitude.
\item \textbf{Perturbation Sound Speed Robustness:} In CAMB, we varied the dark energy rest-frame sound speed across $c_s^2 \in [0.01, 1.00]$. The first acoustic peak remains identically locked at $C_{220}^{\mathrm{TT}} = 5882.66\,\mu\mathrm{K}^2$ (fractional shift $< 10^{-6}\%$), and linear matter clustering is strictly stable at $\sigma_8 = 0.7596$ (PPF mode: $\sigma_8 = 0.7698$). No clustering instabilities or gradient blowups occur.
\item \textbf{Late-Time ISW Spectrum:} Because dark energy evolution is confined to late times ($z < 1$), the only observable CMB perturbation signature is in the low-$\ell$ Integrated Sachs-Wolfe (ISW) spectrum: $+6.11\%$ at $\ell=2$, $+5.63\%$ at $\ell=3$, and $+2.47\%$ at $\ell=10$, while matching $\Lambda$CDM within $<0.8\%$ at $\ell > 100$.
\end{enumerate}

\section{Pre-Registered DESI DR2 Confrontation}
The recent DESI DR2 release (arXiv:2503.14738)~\cite{desi2025} expanded the survey to three years of observations, elevating the dynamical dark energy significance to $3.1\sigma$ (CMB only) and $4.2\sigma$ (CMB + SNe). The observed posterior contours remain firmly anchored in the non-phantom quadrant $w_0 > -1, w_a < 0$, precisely where Horizon-Locked Causal Memory makes its parameter-free prediction.

Future releases from the Roman Space Telescope and the Vera C. Rubin Observatory (LSST) will conclusively test the pre-registered crossover prediction ($E(z) > E_{\Lambda\mathrm{CDM}}$ at $z < 1.2$, $E(z) < E_{\Lambda\mathrm{CDM}}$ at $z > 1.4$).

\section{Conclusion}
Horizon-Locked Causal Memory provides an autonomous, physically grounded resolution to the DESI dynamical dark energy puzzle. By linking past-light-cone causal memory to the cosmic expansion rate $\tau(t) = 1/H(t)$, it derives $(w_0, w_a) = (-0.888, -0.131)$ with zero free parameters, guarantees $w(z) \ge -1$ at all times, eliminates phantom ghosts, and achieves superior statistical significance across all four major contemporary cosmological surveys.

\begin{thebibliography}{99}\footnotesize\setlength{\itemsep}{0pt}
\bibitem{desi2024} DESI Collaboration, arXiv:2404.03002 (2024).
\bibitem{desi2025} DESI Collaboration, arXiv:2503.14738 (2025).
\bibitem{brout2022} D.~Brout et al., Astrophys.\ J.\ \textbf{938}, 110 (2022).
\bibitem{popovic2025} B.~Popovic et al., arXiv:2511.07517 (2025).
\bibitem{planck2018} Planck Collaboration, Astron.\ Astrophys.\ \textbf{641}, A6 (2020).
\bibitem{carroll2003} S.~M.~Carroll, M.~Hoffman and M.~Trodden, Phys.\ Rev.\ D \textbf{68}, 023509 (2003).
\bibitem{hendrixx2026} T.~Hendrixx, \emph{Quantum Relational Causality}, Zenodo preprint (2026).
\bibitem{chen2019} L.~Chen, Q.-G.~Huang and K.~Wang, JCAP \textbf{1902}, 028 (2019).
\bibitem{lewis2000} A.~Lewis, A.~Challinor and A.~Lasenby, Astrophys.\ J.\ \textbf{538}, 473 (2000).
\end{thebibliography}

\end{document}
